\chapter[Alineación con los ODS]{Alineación con los Objetivos de Desarrollo Sostenible}
\label{cap:ods}

El presente Trabajo de Fin de Grado, aunque de naturaleza fundamentalmente técnica y metodológica, se alinea de forma directa e indirecta con varios de los Objetivos de Desarrollo Sostenible (ODS) propuestos por las Naciones Unidas en la Agenda 2030. A continuación, se detalla la contribución del proyecto a los ODS más relevantes.

La tabla \ref{tab:ods} presenta un resumen del grado de relación del trabajo con cada uno de los 17 ODS.

\begin{table}[htbp]
    \centering
    \caption{Grado de relación del trabajo con los Objetivos de Desarrollo Sostenible (ODS).}
    \label{tab:ods}
    \begin{tabular}{|l|c|c|c|c|}
    \hline
    \textbf{Objetivo de Desarrollo Sostenible} & \textbf{Alto} & \textbf{Medio} & \textbf{Bajo} & \textbf{No Procede} \\
    \hline
    ODS 3: Salud y bienestar & \textbf{X} & & & \\
    \hline
    ODS 9: Industria, innovación e infraestructura & \textbf{X} & & & \\
    \hline
    ODS 4: Educación de calidad & & \textbf{X} & & \\
    \hline
    ODS 17: Alianzas para lograr los objetivos & & \textbf{X} & & \\
    \hline
    \end{tabular}
\end{table}

\section{ODS 3: Salud y Bienestar}
\label{sec:ods3}

La contribución más directa y significativa de este proyecto se enmarca en el \textbf{ODS 3: Garantizar una vida sana y promover el bienestar para todos en todas las edades}.

Este trabajo aborda directamente la meta 3.4, que busca reducir la mortalidad por enfermedades no transmisibles, entre las que el cáncer es una de las principales causas. Al mejorar la precisión cuantitativa de la imagen PET con $^{86}\text{Y}$, se impacta directamente en el avance de la \textbf{teranosis}, una de las fronteras de la medicina oncológica personalizada.

Una dosimetría precisa, posibilitada por una imagen PET de alta calidad, permite:
\begin{itemize}
    \item \textbf{Personalizar el Tratamiento:} Ajustar la dosis de radioterapia con $^{90}\text{Y}$ a la biología específica de cada paciente y cada tumor, maximizando la eficacia terapéutica.
    \item \textbf{Aumentar la Seguridad:} Evitar la toxicidad en órganos sanos al predecir con mayor exactitud dónde se acumulará el radiofármaco, mejorando la calidad de vida del paciente.
\end{itemize}

Además, la prueba de concepto de \textbf{imagen multiplexada} (mPET) tiene el potencial de revolucionar el diagnóstico al permitir la visualización simultánea de múltiples procesos biológicos (ej. metabolismo y expresión de receptores). Esto conduciría a una caracterización tumoral mucho más completa en una sola sesión, acelerando el diagnóstico y la toma de decisiones clínicas.

\section{ODS 9: Industria, Innovación e Infraestructura}
\label{sec:ods9}

El proyecto contribuye de forma notable al \textbf{ODS 9: Construir infraestructuras resilientes, promover la industrialización inclusiva y sostenible y fomentar la innovación}.

La colaboración con \textbf{Bruker}, una empresa líder en instrumentación científica, sitúa este TFG en la intersección entre la investigación académica y la innovación industrial. El trabajo se alinea con la meta 9.5, que busca aumentar la investigación científica y mejorar la capacidad tecnológica de los sectores industriales.

Las aportaciones concretas incluyen:
\begin{itemize}
    \item \textbf{Mejora de Tecnología Existente:} La optimización de los parámetros de adquisición para isótopos complejos como el $^{86}\text{Y}$ añade valor a los escáneres PET preclínicos existentes, expandiendo sus capacidades sin necesidad de modificar el hardware.
    \item \textbf{Innovación Algorítmica:} El desarrollo y validación de software de post-procesado para la separación de isótopos (mPET) representa una innovación tecnológica que puede ser transferida a futuros productos comerciales.
    \item \textbf{Infraestructura Científica:} El uso del clúster de supercomputación de la UPV para realizar simulaciones masivas refuerza la importancia de las infraestructuras científicas robustas para el avance de la I+D+i.
\end{itemize}

\section{Otras Contribuciones}
\label{sec:ods_otras}

Adicionalmente, el proyecto tiene una relación de grado medio con los siguientes objetivos:

\begin{itemize}
    \item \textbf{ODS 4: Educación de calidad:} Este TFG, como proyecto final de una titulación de ingeniería, es en sí mismo un producto de la educación superior de calidad. Además, la memoria y el código fuente (si se publican) contribuyen a la diseminación del conocimiento científico y técnico.
    \item \textbf{ODS 17: Alianzas para lograr los objetivos:} La colaboración explícita entre la Universidad (UPV) y la Industria (Bruker) es un ejemplo perfecto de las alianzas necesarias para transferir el conocimiento académico a soluciones con impacto real en la sociedad.
\end{itemize}